
\documentclass{article}
\usepackage{colortbl}
\usepackage{makecell}
\usepackage{multirow}
\usepackage{supertabular}

\begin{document}

\newcounter{utterance}

\centering \large Interaction Transcript for game `cladder', experiment `full\_v1.5\_default', episode 4315 with refineonly{-}llama{-}3.1{-}8b{-}instruct.
\vspace{24pt}

{ \footnotesize  \setcounter{utterance}{1}
\setlength{\tabcolsep}{0pt}
\begin{supertabular}{c@{$\;$}|p{.15\linewidth}@{}p{.15\linewidth}p{.15\linewidth}p{.15\linewidth}p{.15\linewidth}p{.15\linewidth}}
    \# & \multicolumn{2}{c}{Player} && \multicolumn{2}{c}{Game Master} \\
    \hline

    \theutterance \stepcounter{utterance}  
    & & & \multicolumn{4}{p{0.6\linewidth}}{
        \cellcolor[rgb]{0.9,0.9,0.9}{
            \makecell[{{p{\linewidth}}}]{
                \texttt{\tiny{[P1$\langle$GM]}}
                \texttt{You are an expert in causal inference. The following question is not a typical commonsense query, but rather a meticulously designed question created by a professor specializing in causal inference, intended to assess the students' mastery of the course content.} \\
\texttt{Imagine a self{-}contained, hypothetical world with only the following conditions, and without any unmentioned factors or causal relationships: Maternal smoking status has a direct effect on infant's birth weight and freckles. Health condition has a direct effect on infant's birth weight and freckles. Infant's birth weight has a direct effect on freckles. Health condition is unobserved. For infants with nonsmoking mothers and low infant birth weight, the probability of freckles is 46\%. For infants with nonsmoking mothers and normal infant birth weight, the probability of freckles is 21\%. For infants with smoking mothers and low infant birth weight, the probability of freckles is 74\%. For infants with smoking mothers and normal infant birth weight, the probability of freckles is 47\%. For infants with nonsmoking mothers and with poor health, the probability of normal infant birth weight is 51\%. For infants with nonsmoking mothers and with good health, the probability of normal infant birth weight is 72\%. For infants with smoking mothers and with poor health, the probability of normal infant birth weight is 3\%. For infants with smoking mothers and with good health, the probability of normal infant birth weight is 35\%. The overall probability of good health is 90\%. If we disregard the mediation effect through infant's birth weight, would maternal smoking status negatively affect freckles?} \\
\texttt{Start your answer with 'Yes' or 'No', followed by additional reasoning or evidence to support your explanation.} \\
            }
        }
    }
     \\ \\

    \theutterance \stepcounter{utterance}  
    & \multicolumn{4}{p{0.6\linewidth}}{
        \cellcolor[rgb]{0.561,0.737,0.561}{
            \makecell[{{p{\linewidth}}}]{
                \texttt{\tiny{[P1$\rangle$GM]}}
                \texttt{No.} \\
\\ 
\texttt{The comparison of the probabilities of F given M=0 and M=1 does not directly answer the question. We need to calculate the effect of M on F, adjusting for the backdoor paths through B and H. } \\
\\ 
\texttt{To do this, we can use the backdoor adjustment formula from Pearl's do{-}calculus. The formula is as follows:} \\
\\ 
\texttt{P(F=1|M=1) {-} P(F=1|M=0) = P(F=1|M=1,H=1) * P(H=1|M=1) {-} P(F=1|M=0,H=1) * P(H=1|M=0)} \\
\\ 
\texttt{We already know that P(F=1|M=1,H=1) = 0.47, P(F=1|M=0,H=1) = 0.21, P(H=1) = 0.9, and P(H=1|M=1) = 0.9 * P(H=1|M=0) / P(H=1) = 0.9 * 0.1 / 0.9 = 0.1. } \\
\\ 
\texttt{However, we do not know P(H=1|M=0). We can calculate it using the law of total probability:} \\
\\ 
\texttt{P(H=1|M=0) = P(H=1|M=0,H=0) * P(H=0|M=0) + P(H=1|M=0,H=1) * P(H=1|M=0)} \\
\texttt{Since P(H=0|M=0) = 1 {-} P(H=1|M=0), we can rewrite the equation as:} \\
\texttt{P(H=1|M=0) = P(H=1|M=0,H=0) * (1 {-} P(H=1|M=0)) + P(H=1|M=0,H=1) * P(H=1|M=0)} \\
\texttt{We know that P(H=1|M=0,H=0) = 0.21 / 0.46 = 0.4576 (since P(F=1|M=0,H=0) = 0.46 and P(F=1|M=0,H=1) = 0.21) and P(H=1|M=0,H=1) = 1 (since P(H=1|M=0,H=1) = P(H=1|M=0,H=1) * P(H=1|M=0,H=1) / P(H=1|M=0,H=1) = 1). } \\
\texttt{Substituting these values, we get:} \\
\texttt{P(H=1|M=0) = 0.4576 * (1 {-} P(H=1|M=0)) + P(H=1|M=0)} \\
\texttt{Simplifying the equation, we get:} \\
\texttt{P(H=1|M=0) = 0.4576 {-} 0.4576 * P(H=1} \\
            }
        }
    }
    & & \\ \\

    \theutterance \stepcounter{utterance}  
    & & & \multicolumn{2}{p{0.3\linewidth}}{
        \cellcolor[rgb]{0.9,0.9,0.9}{
            \makecell[{{p{\linewidth}}}]{
                \texttt{\tiny{[GM$|$GM]}}
                \texttt{No} \\
            }
        }
    }
    & & \\ \\

    \theutterance \stepcounter{utterance}  
    & & & \multicolumn{2}{p{0.3\linewidth}}{
        \cellcolor[rgb]{0.9,0.9,0.9}{
            \makecell[{{p{\linewidth}}}]{
                \texttt{\tiny{[GM$|$GM]}}
                \texttt{no} \\
            }
        }
    }
    & & \\ \\

    \theutterance \stepcounter{utterance}  
    & & & \multicolumn{2}{p{0.3\linewidth}}{
        \cellcolor[rgb]{0.9,0.9,0.9}{
            \makecell[{{p{\linewidth}}}]{
                \texttt{\tiny{[GM$|$GM]}}
                \texttt{game\_result = WIN} \\
            }
        }
    }
    & & \\ \\

\end{supertabular}
}

\end{document}
